\section*{Capítulo \ref{ch:b2:catalytic-cycles} --- Catálise organometálica: etapas elementares e ciclos}

\begin{solution}{exo:b2:catalytic-cycles:1}
Antes: Ir(I), $d^8$, $8 + 4 \times 2 = 16$ elétrons. Depois: Ir(III), $d^6$, $6 + 6 \times 2 = 18$
elétrons, octaédrico.
\end{solution}

\begin{solution}{exo:b2:catalytic-cycles:2}
A primeira é uma \omterm{def:b2:catalytic-cycles:oxidative-addition}{eliminação redutiva} (de Pt(IV) a Pt(II), com etano formado a
partir de duas metilas). A segunda é uma \omterm{def:b2:catalytic-cycles:ligand-exchange}{troca de ligante} (uma fosfina
substituída pelo eteno) seguida de uma \omterm{def:b2:catalytic-cycles:insertion}{inserção migratória} do eteno na
ligação Pd--Ph.
\end{solution}

\begin{solution}{exo:b2:catalytic-cycles:3}
TON $= 15.0/0.020 = 750$; TOF $= 750/3.0 = \qty{250}{h^{-1}}$.
\end{solution}

\begin{solution}{exo:b2:catalytic-cycles:4}
\begin{align*}
  &\mathrm{PdL_2 + ArBr \to ArPdBrL_2}\\
  &\mathrm{ArPdBrL_2 + Ar'B(OH)_3^- \to ArPdAr'L_2 + Br^- + B(OH)_3}\\
  &\mathrm{ArPdAr'L_2 \to Ar{-}Ar' + PdL_2}
\end{align*}
Soma: $\mathrm{ArBr + Ar'B(OH)_3^- \to Ar{-}Ar' + Br^- + B(OH)_3}$; todas as espécies de
paládio se cancelam.
\end{solution}

\begin{solution}{exo:b2:catalytic-cycles:5}
Etila e isopropila (elas têm hidrogênios $\beta$). A metila não tem carbono $\beta$;
o carbono $\beta$ da neopentila não tem hidrogênio; o carbono $\beta$ da benzila é um carbono
\omterm{def:b2:huckel:aromatic}{aromático} do anel, sem hidrogênio.
\end{solution}

\begin{solution}{exo:b2:catalytic-cycles:6}
O (\textit{E})-3-fenilpropenoato de metila, \ce{Ph-CH=CH-COOCH3}, com a fenila
adicionada ao carbono terminal; produto trans.
\end{solution}

\begin{solution}{exo:b2:catalytic-cycles:7}
Butanal \ce{CH3CH2CH2CHO} (linear: o metal se adiciona ao carbono terminal, o
hidreto ao interno) e 2-metilpropanal \ce{(CH3)2CHCHO} (ramificado: o metal no
carbono interno).
\end{solution}

\begin{solution}{exo:b2:catalytic-cycles:8}
O \ce{[RhCl(PPh3)3]} tem 16 elétrons: adicionar \ce{H2} (+2) daria 18 com seis
posições de coordenação, cheio demais com três fosfinas volumosas. O
\ce{[RhCl(PPh3)2]} tem 14 e um sítio livre: a \omterm{def:b2:catalytic-cycles:oxidative-addition}{adição oxidativa} dá facilmente
um complexo de 16 elétrons.
\end{solution}

\begin{solution}{exo:b2:catalytic-cycles:9}
O ácido borônico é um nucleófilo fraco; a base adiciona hidróxido (ou alcóxido)
ao boro, dando um borato \ce{ArB(OH)3-}, cujo grupo arila é mais rico em
elétrons e se transfere para o paládio.
\end{solution}

\begin{solution}{exo:b2:catalytic-cycles:10}
Um complexo de 20 elétrons não é razoável. O complexo precisa primeiro perder um
ligante (14 e), depois adicionar \ce{H2} (16 e), ligar o alceno (18 e),
inseri-lo em uma ligação M--H (16 e) e eliminar o alcano (14 e), antes de
retomar o ligante ou recomeçar.
\end{solution}

\begin{solution}{exo:b2:catalytic-cycles:11}
Velocidade inicial $n_{\max}/\tau = \qty{0.25}{mmol/min}$; TOF inicial $= 0.25/0.010 =
\qty{25}{min^{-1}} = \qty{1.5e3}{h^{-1}}$. Final TON $= 10.0/0.010 = 1.0 \times 10^3$.
\end{solution}

\begin{solution}{exo:b2:catalytic-cycles:12}
4-Bromoanisol com ácido fenilborônico (ou bromobenzeno com ácido
4-metoxifenilborônico), um catalisador fosfina-paládio(0) e uma base. Ciclo:
\omterm{def:b2:catalytic-cycles:oxidative-addition}{adição oxidativa} do brometo de arila, \omterm{def:b2:catalytic-cycles:transmetalation}{transmetalação} a partir do borato,
\omterm{def:b2:catalytic-cycles:oxidative-addition}{eliminação redutiva} do 4-metoxibifenila.
\end{solution}

\begin{solution}{pb:b2:catalytic-cycles:1}
\textbf{1.} Pd(II), $d^8$.
\textbf{2.} Pd(0), $d^{10}$, $10 + 2 \times 2 = 14$ elétrons.
\textbf{3.} Com 14 elétrons, ele pode aceitar mais dois pares: tem \omterm{def:b2:catalytic-cycles:unsaturated}{sítios vacantes}.
\textbf{4.} Os complexos fosfina-paládio(0) são oxidados pelo ar, e as fosfinas
são oxidadas a óxidos de fosfina.
\textbf{5.} Ela reduz o paládio(II) a paládio(0) (sendo ela mesma oxidada) e
liga o paládio(0), mantendo-o em solução.
\textbf{6.} O \omterm{def:b2:catalytic-cycles:cycle}{pré-catalisador}.
\textbf{7.} \ce{Pd(PPh3)2 + CH3C6H4Br -> [Pd(C6H4CH3)Br(PPh3)2]}: Pd(II), $d^8$, $8 + 4 \times 2 =
16$ elétrons.
\textbf{8.} Ele o transforma no borato \ce{PhB(OH)3-}, que transfere seu grupo fenila.
\textbf{9.} \ce{[Pd(Ar)Br(PPh3)2] + PhB(OH)3- -> [Pd(Ar)(Ph)(PPh3)2] + Br- + B(OH)3}.
\textbf{10.} A \omterm{def:b2:catalytic-cycles:oxidative-addition}{eliminação redutiva} une dois ligantes cis; arilas em trans precisam
primeiro se isomerizar.
\textbf{11.} \ce{[Pd(Ar)(Ph)(PPh3)2] -> Ar-Ph + Pd(PPh3)2}: paládio(0), 14 elétrons, pronto
para outra volta.
\textbf{12.} \ce{CH3C6H4Br + PhB(OH)3- -> CH3C6H4-Ph + Br- + B(OH)3}.
\textbf{13.} O 4-bromotolueno (\qty{10.0}{mmol}, contra \qty{12.0}{mmol} de ácido borônico).
\textbf{14.} $1.51/168.24 = \qty{8.98e-3}{mol}$, \qty{8.98}{mmol}.
\textbf{15.} $8.98/10.0 = 89.8~\%$.
\textbf{16.} $0.050/10.0 = 0.50$~mol~\%.
\textbf{17.} $8.98/0.050 = 180$.
\textbf{18.} $180/4.0 = \qty{45}{h^{-1}}$.
\textbf{19.} Dois grupos fenila transferidos para o mesmo paládio
(homoacoplamento do ácido borônico, favorecido por traços de oxigênio) e
depois eliminados juntos.
\textbf{20.} Os grupos arila não têm hidrogênio $\beta$ capaz de alcançar o metal.
\textbf{21.} A \omterm{def:b2:catalytic-cycles:oxidative-addition}{adição oxidativa}: a ligação \ce{C-Cl} é mais forte que \ce{C-Br}.
\textbf{22.} Ele termina como partículas de paládio(0) ou sais nos resíduos; é
caro e tóxico, e é recuperado.
\textbf{23.} Parte do ácido borônico se perde em reações secundárias
(homoacoplamento, perda do boro do anel); o excesso mantém o brometo como
reagente limitante.
\textbf{24.} O paládio realizou $\boldsymbol{\approx \num{1.8e2}}$ rotações.
\end{solution}
